\section{Nef b-line decompositions}\label{sec:supply}

\begin{proposition}\label{prop:expression}
Let $(X,B)$ be a globally strongly $\Q$-factorial compact Kähler
fourfold pair, with $B$ effective and rational, $(X,B)$ klt, and
$D=K_X+B$ rational Cartier and analytically pseudo-effective. Suppose
there is a projective log resolution $p:W\to X$, with $W$ smooth and
compact Kähler, such that $W$ is non-uniruled or Moishezon. Then there
exist a smooth compact Kähler manifold $U$, a
projective modification $r:U\to X$, an analytically nef rational
holomorphic line bundle $M_U$, and an effective rational divisor $N$ on
$X$ such that, for the rational b-line bundle $\bM$ determined by $M_U$,
\begin{equation}\label{eq:supply-expression}
 D\simq N+M_X.
\end{equation}
The equivalence is an identity of actual rational reflexive sheaves.
\end{proposition}

\begin{proof}
Choose a projective log resolution $p:W\to X$ as in the statement.
We first describe how to transfer a nef adjoint on a minimal model reached
by a finite program from $W$ to a rational b-line summand over $X$.

Suppose that a finite program on $W$ ends at a model
$Y$, and that $H_W$ and $H_Y$ are its initial and final actual rational
adjoints.  Both morphisms defining each step are projective.  The main
components of the successive fiber products of their graphs, followed
by normalization and resolution, therefore give a diagram
\[
 \begin{array}{ccc}
 &U&\\[-2pt]
 a\swarrow&&\searrow q\\[-2pt]
 W&&Y
 \end{array}
\]
whose two arrows are projective modifications.  The construction uses
only the finitely many steps of this program: projectivity follows
successively by base change and composition.  In particular $r=pa$ is
projective over the original $X$, and $U$ is compact Kähler.

The stepwise comparisons in Lemma~\ref{lem:negativity}, pulled back to
this common model, give
\begin{equation}\label{eq:supply-common-comparison}
 a^*H_W\simq q^*H_Y+G,\qquad G\geq0.
\end{equation}
Set $M_U=q^*H_Y$.  When $H_Y$ is nef, so is $M_U$: pull back metrics with
arbitrarily small negative curvature and compare the pulled-back fixed
form with a Kähler form on $U$.  In the projective case $U$ is itself
projective, and the pullback of a nef rational bundle is analytically
nef by the projective metric characterization.  For a common multiple
$m$, the trace of $M_U$ is
represented by the rank-one reflexive sheaf
\[
 \bigl(r_*\cO_U(mM_U)\bigr)^{**}.
\]
Coherence follows from properness, and global strong $\Q$-factoriality
on $X$ makes a further reflexive power invertible.  By
Lemma~\ref{lem:trace}, tracing \eqref{eq:supply-common-comparison} gives
\begin{equation}\label{eq:supply-trace-comparison}
 \operatorname{Tr}_r(a^*H_W)\simq M_X+r_*G.
\end{equation}
Indeed this is the usual tensor identity at each codimension-one point
of $X$, where $r$ is an isomorphism, and both sides extend reflexively.
It retains the actual identifications of the line bundles.  In
particular, no global meromorphic frame of either canonical bundle or
of $M_U$ is chosen.

\medskip\noindent
\textit{Non-uniruled resolution.}
By \cite[Theorem~1.1]{Ou2025}, $K_W$ is analytically pseudo-effective.
Starting with the smooth empty-boundary pair $(W,0)$, apply
Lemma~\ref{lem:absolute-support} and Proposition~\ref{prop:ordinary-step}
whenever the current canonical bundle is not nef.
Lemma~\ref{lem:pe-step} preserves its
pseudo-effectivity, and Lemma~\ref{lem:mori-not-pe} excludes a
fiber-type contraction at every stage.  Terminality persists by
Lemma~\ref{lem:negativity}.  Corollary~\ref{prop:terminal-four} rules
out an infinite sequence.  Continuation from every non-nef model thus
gives a finite program ending at a terminal compact Kähler $Y$ with
analytically nef actual rational canonical bundle $K_Y$.

Apply \eqref{eq:supply-common-comparison} with $H_W=K_W$ and $H_Y=K_Y$.
The reflexive trace of $a^*K_W$ on $X$ is $K_X$: over the common
codimension-one open set the canonical sheaves are canonically
identified, and normality determines their reflexive extensions.
Equation~\eqref{eq:supply-trace-comparison} therefore yields
\[
 K_X+B\simq M_X+(r_*G+B).
\]
Take $N=r_*G+B$.  This argument uses the trace of $K_W$; it does not
require the discrepancies comparing $K_W$ with $p^*D$ to be
nonnegative.  The terminal fourfold result used here was established
independently of Theorem~\ref{thm:expression-termination}.

\medskip\noindent
\textit{Uniruled Moishezon resolution.}
The smooth Kähler manifold $W$ is then projective.  By
Lemma~\ref{lem:supply-log-resolution}, choose an effective rational SNC
klt boundary $\Gamma$ satisfying
\begin{equation}\label{eq:supply-projective-start}
 H_W:=K_W+\Gamma\simq p^*D+E_0,
 \qquad E_0\geq0\quad\text{and }E_0\text{ is }p\text{-exceptional}.
\end{equation}
This adjoint is analytically pseudo-effective, hence numerically
pseudo-effective in the projective sense by
\cite[Proposition~1.2]{BDPP2013}.

Run the ordinary projective klt MMP on $(W,\Gamma)$.  Divisorial steps
are finite, and every remaining flip sequence terminates by
\cite[Theorem~1.1]{CT2023}, applied to a pseudo-effective NQC lc
fourfold with zero nef part.  These are projective varieties over a
point, as required by that theorem's conventions.  Pseudo-effectivity
persists by Lemma~\ref{lem:pe-step}, so Lemma~\ref{lem:mori-not-pe}
excludes a fiber-type negative contraction.  The resulting model $Y$ has nef
rational adjoint $H_Y=K_Y+\Gamma_Y$.

Form \eqref{eq:supply-common-comparison} for this finite program.
Since $E_0$ is exceptional over $X$, the reflexive trace of
$a^*H_W$ in \eqref{eq:supply-projective-start} is $D$.  Thus
\eqref{eq:supply-trace-comparison} gives
\[
 D\simq M_X+r_*G.
\]
Take $N=r_*G$.  The projective fourfold theorem has supplied a nef
adjoint; fourfold abundance is not used.
\end{proof}
